\section*{Hoofdstuk \ref{ch:b2:atmospheric-thermodynamics} --- Thermodynamica van de atmosfeer}

\begin{solution}{exo:b2:atmospheric-thermodynamics:1}
(a) $e = e_s(\qty{278.15}{K}) = \qty{8.73}{hPa}$ en $e_s(\qty{283.15}{K}) = \qty{12.32}{hPa}$.
(b) $8.73/12.32 = \qty{70.9}{\percent}$. (c) $r = 0.622 \times 8.73/(850 - 8.73) =
\qty{6.45}{g/kg}$ en $q = r/(1 + r) = \qty{6.41}{g/kg}$.
\end{solution}

\begin{solution}{exo:b2:atmospheric-thermodynamics:2}
(a) $\mathrm{d}\ln e_s/\mathrm{d}T = 5419/T^2$: \qty{7.26}{\percent} per kelvin bij
\qty{273.15}{K}, \qty{6.31}{\percent} bij \qty{293.15}{K}, \qty{5.90}{\percent} bij
\qty{303.15}{K}. (b) $e_s(\qty{20}{\celsius}) = \qty{23.66}{hPa}$ en $e_s(\qty{25}{\celsius})
= \qty{32.26}{hPa}$; verhouding 1.363, tegen $1.063^5 = 1.357$. De snelheid is een afgeleide van
$\ln e_s$, dus over \qty{5}{K} is de factor $\exp(5 \times 0.0620) = 1.363$, met 0.0620 de
snelheid in het midden, \qty{22.5}{\celsius}; vijf keer \qty{6.3}{\percent} samengesteld komt dicht in de buurt,
maar is niet exact. (c) $23.66/23.39 = 1.012$: de formule is \qty{1.2}{\percent} te hoog, omdat
ze de latente warmte van het tripelpunt aanhoudt, terwijl de latente warmte afneemt naarmate de temperatuur
stijgt; de helling van de formule is te steil, en haar fout groeit met de afstand tot
\qty{0}{\celsius}.
\end{solution}

\begin{solution}{exo:b2:atmospheric-thermodynamics:3}
(a) $125 \times (30 - 14) = \qty{2000}{m}$. (b) $30 - 9.8 \times 2.0 = \qty{10.4}{\celsius}$.
(c) $14 - 1.8 \times 2.0 = \qty{10.4}{\celsius}$: temperatuur en dauwpunt ontmoeten elkaar, zoals ze moeten op
het verzadigingsniveau; de regel is deze berekening. (d) $125 \times (33 - 14) = \qty{2375}{m}$:
de bases stijgen naarmate de middag de lucht bij de grond opwarmt.
\end{solution}

\begin{solution}{exo:b2:atmospheric-thermodynamics:4}
(a) De dauwpuntskromme ligt op de temperatuurkromme van de grond tot ongeveer \qty{908}{hPa}:
de onderste halve kilometer is verzadigd, of bijna. (b) Tussen ongeveer 908 en \qty{889}{hPa}
buigt de temperatuurkromme naar rechts, naar warmere waarden, terwijl de druk afneemt: de
temperatuur stijgt met de hoogte. (c) LCL bij ongeveer \qty{964}{hPa}, enkele tientallen meters hoger; LFC bij
ongeveer \qty{730}{hPa}; EL bij ongeveer \qty{235}{hPa}. (d) Aan de top van de inversie ligt het dauwpunt
\qty{11}{K} onder de temperatuur, en \qty{20}{K} eronder bij \qty{850}{hPa}: warme, zeer droge
lucht ligt boven koele, vochtige lucht. Haar relatieve vochtigheid is laag, dus ze kreeg haar
warmte niet door condensatie; lucht die daalt, warmt op door compressie met de droogadiabatische gradiënt terwijl
ze haar \omterm{def:b2:atmospheric-thermodynamics:moist-air}{mengverhouding} behoudt, wat zo'n laag zou geven.
\end{solution}

\begin{solution}{exo:b2:atmospheric-thermodynamics:5}
(a) $T_v = 303.15 \times (1 + 0.020/0.622)/1.020 = \qty{306.76}{K}$; de benadering geeft
$303.15 \times (1 + 0.61 \times 0.020) = \qty{306.85}{K}$. (b) $\rho = p/(R_dT_v) = 10^5/(287.0
\times 306.76) = \qty{1.1358}{kg/m^3}$, tegen $10^5/(287.0 \times 303.15) =
\qty{1.1494}{kg/m^3}$ voor droge lucht: \qty{1.2}{\percent} lichter. (c) Met $T_v - T =
\qty{3.6}{K}$. (d) Een stijgend pakketje hoeft alleen lichter te zijn dan zijn omgeving: boven warme zeeën
voegt het vocht van de lucht een opwaartse kracht toe die gelijkstaat aan enkele kelvin, evenveel als een flink
temperatuuroverschot.
\end{solution}

\begin{solution}{exo:b2:atmospheric-thermodynamics:6}
(a) $e_s$ = 6.11, 23.66 en \qty{43.54}{hPa}; $r_s = 0.622\,e_s/(1000 - e_s)$ = 3.82, 15.07 en
\qty{28.31}{g/kg}; $\Gamma_s$ = 6.48, 4.20 en \qty{3.44}{K/km}. (b) Het droge pakketje koelt
\qty{9.8}{K} af, het verzadigde \qty{4.2}{K}: het verzadigde pakketje is \qty{5.6}{K} warmer.
(c) Hogerop is de lucht kouder, is $r_s$ kleiner, blijft er weinig damp over om te condenseren, en
keert $\Gamma_s$ terug naar $\Gamma_d$.
\end{solution}

\begin{solution}{exo:b2:atmospheric-thermodynamics:7}
(a) $e = e_s(\qty{291.65}{K}) = \qty{21.52}{hPa}$, $r = 0.622 \times 21.52/(971.1 - 21.52) =
\qty{14.1}{g/kg}$; $\theta = 292.15 \times (1000/971.1)^{0.2857} = \qty{294.6}{K}$; $\theta_e =
294.6 \exp[2.501 \times 10^6 \times 0.0141/(1004.5 \times 291.5)] = 294.6 \times 1.128 =
\qty{332.3}{K}$. (b) $e_s(\qty{264.25}{K}) = \qty{3.13}{hPa}$, $r_s = 0.622 \times 3.13/(500 -
3.13) = \qty{3.92}{g/kg}$; $\theta = 264.25 \times 2^{0.2857} = \qty{322.1}{K}$; het pakketje is
verzadigd, dus $T_{\text{LCL}} = T$ en $\theta_e = 322.1 \exp[2.501 \times 10^6 \times
0.00392/(1004.5 \times 264.25)] = \qty{334.3}{K}$. (c) $\theta$ steeg met \qty{27.5}{K}, opgewarmd
door de latente warmte van de \qty{10.2}{g/kg} damp die condenseerde; $\theta_e$ veranderde met
\qty{0.6}{\percent}, binnen de nauwkeurigheid van zijn benaderende formule: hij blijft behouden.
\end{solution}

\begin{solution}{exo:b2:atmospheric-thermodynamics:8}
(a) $\text{CAPE} = 287.0 \times 4 \times \ln(700/300) = 287.0 \times 4 \times 0.847 =
\qty{973}{J/kg}$. (b) $\sqrt{2 \times 973} = \qty{44.1}{m/s}$. (c) Tussen 1 en
\qty{1500}{J/kg}: positief, maar niet groot. (d) Een band met constante horizontale breedte tussen twee
krommen; volgens \cref{prop:b2:atmospheric-thermodynamics:areas} is haar oppervlakte evenredig met $\int
\Delta T\,\mathrm{d}\ln p = 4 \ln(700/300)$, wat de helling van de isothermen ook is.
\end{solution}

\begin{solution}{exo:b2:atmospheric-thermodynamics:9}
(a) $D_c = \sqrt{3 \times 0.61 \times 80^3/2.1}\,\unit{nm} = \qty{668}{nm} =
\qty{0.668}{\micro m}$; $s_c = 2A/(3D_c) = 2 \times 2.1/(3 \times 668) = 0.00210 =
\qty{0.210}{\percent}$. (b) $s_c = \sqrt{4A^3/(27\kappa D_d^3)}$ oplossen naar $D_d$: $D_d =
\bigl(4A^3/(27\kappa s_c^2)\bigr)^{1/3}$, \qty{63}{nm} bij \qty{0.3}{\percent} en \qty{131}{nm} bij
\qty{0.1}{\percent}. (c) Elk deeltje tussen 63 en \qty{131}{nm} wordt nu ook geactiveerd: meer
druppels, en, bij hetzelfde vloeibare water, kleinere.
\end{solution}

\begin{solution}{exo:b2:atmospheric-thermodynamics:10}
(a) $w = \sqrt{2 \times 130} = \qty{16.1}{m/s}$, onmogelijk voor lucht in rust bij de grond.
(b) De potentiële temperatuur van de top van de inversie is $293.65 \times (1000/889.1)^{0.2857} =
\qty{303.7}{K}$; een oppervlaktepakketje met deze $\theta$ heeft $T = 303.7 \times (971.1/1000)^{0.2857}
= \qty{301.2}{K}$, \qty{28.0}{\celsius}. (c) $125 \times (28.0 - 18.5) = \qty{1190}{m}$ boven de
grond, boven de oude top van de inversie. (d) Een front, dat de lucht langs zijn helling optilt; de trage stijging
van de lucht vóór een storing in de hoogte, die het deksel zelf afkoelt.
\end{solution}

\begin{solution}{exo:b2:atmospheric-thermodynamics:11}
(a) Wanneer $r_s \to 0$, gaat de teller naar $g$ en de noemer naar $c_p$. (b) Schrijf
$L^2r_s/(R_vT^2) = c_p \times \bigl(Lr_s/(R_dT)\bigr) \times LR_d/(R_vTc_p)$, met
$LR_d/(R_vTc_p) = 5.1$ bij \qty{303}{K}: wanneer $Lr_s/(R_dT) \gg 1$, is ook $L^2r_s/(R_vT^2) \gg c_p$,
en $\Gamma_s \to g\,(Lr_s/(R_dT))/(L^2r_s/(R_vT^2)) = gR_vT/(LR_d) = 9.81 \times 461.5 \times
303/(2.501 \times 10^6 \times 287.0) = \qty{1.91}{K/km}$. (c) Met $T_d = T$ is
$\Gamma_{\text{dauw}} = gR_vT/(LR_d)$, hetzelfde. In die limiet is de latente warmte zo groot dat
het pakketje bijna niets van zijn damp kan condenseren zonder ver boven zijn
verzadigingspunt op te warmen: het behoudt zijn \omterm{def:b2:atmospheric-thermodynamics:moist-air}{mengverhouding} en volgt de lijn van constante \omterm{def:b2:atmospheric-thermodynamics:moist-air}{mengverhouding}, waarlangs de
temperatuur daalt met $\Gamma_{\text{dauw}}$. (d) Bij \qty{30}{\celsius} en \qty{1000}{hPa} is
$Lr_s/(R_dT) = 0.81$, niet veel groter dan 1, en $L^2r_s/(R_vT^2) = 4180$, vier keer $c_p$:
$\Gamma_s = \qty{3.4}{K/km}$, nog ver van \qty{1.9}{K/km}. De limiet wordt niet bereikt.
\end{solution}

\begin{solution}{exo:b2:atmospheric-thermodynamics:12}
(a) $2.868/2.599 = 1.104$ en $1.257/1.032 = 1.218$: \qty{10.4}{\percent} en
\qty{21.8}{\percent}. (b) Druppels worden geactiveerd bij enkele tienden van een procent \omterm{def:b2:atmospheric-thermodynamics:microphysics}{oververzadiging},
tientallen keren minder. (c) De kristallen groeien snel door depositie; ze nemen damp uit de lucht,
die onder verzadiging boven de druppels zakt; de druppels verdampen en hun damp gaat naar
het ijs (het Bergeronproces), tot de wolk volledig bevroren is. (d) Bij \qty{0}{\celsius} zijn de twee
verzadigingsdampdrukken gelijk: ze ontmoeten elkaar in het tripelpunt, en de verhouding gaat
daar dichtbij naar één.
\end{solution}

\begin{solution}{pb:b2:atmospheric-thermodynamics:1}
\textbf{1.} $e = e_s(\qty{291.65}{K}) = \qty{21.5}{hPa}$.
\textbf{2.} $r = 0.622 \times 21.5/(971.1 - 21.5) = \qty{14.1}{g/kg}$.
\textbf{3.} $q = 14.1/1.0141 = \qty{13.9}{g/kg}$.
\textbf{4.} $e_s(\qty{292.15}{K}) = \qty{22.2}{hPa}$; $21.5/22.2 = \qty{97}{\percent}$.
\textbf{5.} $0.61 \times 0.0141 \times 292.15 = \qty{2.5}{K}$.
\textbf{6.} $125 \times 0.5 = \qty{63}{m}$; $p = 971.1 \exp[-9.81 \times 62.5/(287.0 \times 294.7)]
= \qty{964}{hPa}$.
\textbf{7.} $\theta = 292.15 \times (1000/971.1)^{0.2857} = \qty{294.6}{K}$.
\textbf{8.} $16.9 - 17.3 = \qty{-0.4}{K}$: kouder.
\textbf{9.} $(20.5 - 16.4)/(1102 - 924) = 4.1/178$: de temperatuur stijgt met \qty{23}{K/km}, een
gradiënt van \qty{-23}{K/km}. Een inversie: absoluut stabiel, een deksel.
\textbf{10.} $15.4 - 20.5 = \qty{-5.1}{K}$.
\textbf{11.} $e_s(\qty{-1.8}{\celsius}) = \qty{5.4}{hPa}$ en $e_s(\qty{18.6}{\celsius}) =
\qty{21.7}{hPa}$: \qty{25}{\percent}. Warm en droog kan het lucht zijn die daalde en door
compressie opwarmde, of lucht die bovenwinds boven droge grond werd verhit; het is niet de vochtige lucht eronder.
\textbf{12.} $18.6 + 18.4 = \qty{37.0}{K}$ over $5725 - 1489 = \qty{4236}{m}$, \qty{8.73}{K/km}: tussen $\Gamma_s$ en
$\Gamma_d$, voorwaardelijk onstabiel.
\textbf{13.} $-18.4 - (-8.9) = \qty{-9.5}{K}$: extreme instabiliteit ($-8$ tot $-10$).
\textbf{14.} Van de grond omhoog: 0, $-0.4$, $-0.2$, $-5.1$, $-4.9$, $+1.6$, $+9.5$, $+10.0$,
$+5.6$, $+1.5$, $-9.5$~\unit{K}.
\textbf{15.} Tussen 850 ($-4.9$) en \qty{700}{hPa} ($+1.6$), bij de fractie $4.9/6.5 = 0.754$:
$\ln p = \ln 850 + 0.754 \ln(700/850)$, $p = \qty{734}{hPa}$; $z = 1489 + 0.754 \times 1621 =
\qty{2711}{m}$, \qty{2.37}{km} boven de grond.
\textbf{16.} Tussen 250 ($+1.5$) en \qty{200}{hPa} ($-9.5$), bij de fractie $1.5/11.0 =
0.136$: $p = \qty{242.5}{hPa}$ en $z = 10560 + 0.136 \times 1446 = \qty{10757}{m}$, ongeveer
\qty{200}{m} onder de gemelde tropopauze (\qty{10954}{m}).
\textbf{17.} De segmenten (gemiddelde $\Delta T$ maal $\ln$ van de drukverhouding): $0.80 \times 0.0474
= 0.038$; $5.55 \times 0.3365 = 1.867$; $9.75 \times 0.2231 = 2.176$; $7.80 \times 0.2877 =
2.244$; $3.55 \times 0.1823 = 0.647$; $0.75 \times 0.0305 = 0.023$; som \qty{6.995}{K}.
$\text{CAPE} = 287.0 \times 6.995 = \qty{2010}{J/kg}$, op drie cijfers.
\textbf{18.} Tussen 1500 en \qty{2500}{J/kg}: groot.
\textbf{19.} $-0.2 \times 0.0486 - 0.3 \times 0.0189 - 2.65 \times 0.0207 - 5.0 \times 0.0450 -
2.45 \times 0.1467 = -0.655$~\unit{K}; $\text{CIN} = 287.0 \times (-0.655) = \qty{-188}{J/kg}$.
\textbf{20.} De tabel houdt maar elf niveaus aan: daartussen vervangt de trapeziumregel de
krommen door rechte stukken in $\ln p$ en mist ze hun knikken. Het rooster van \qty{1}{hPa} volgt de krommen. De \omterm{def:b2:atmospheric-thermodynamics:virtual}{virtuele temperatuur}
voegt de opwaartse kracht van de waterdamp van het pakketje toe, die overvloediger is in het verzadigde pakketje dan in
de droge lucht eromheen: een grotere CAPE, een kleinere CIN.
\textbf{21.} $\sqrt{2 \times 2010} = \qty{63}{m/s}$.
\textbf{22.} $10757 - 2711 = \qty{8046}{m}$ met $63.4/2 = \qty{31.7}{m/s}$: \qty{254}{s}, ongeveer vier
minuten.
\textbf{23.} Het gewicht van het gecondenseerde water dat de opwaartse stroming meedraagt; de bijmenging van drogere lucht
eromheen (entrainment), die haar afkoelt doordat druppels verdampen.
\textbf{24.} De pakketjestheorie geeft de opwaartse stroming van een onweersbui die door deze lucht wordt gevoed een snelheid van ongeveer
\textbf{\qty[text-series-to-math=true]{63}{m.s^{-1}}}, een bovengrens.
\end{solution}
